\documentclass[../../../include/open-logic-section]{subfiles}

\begin{document}

\olfileid{sth}{ordinals}{opps}
\olsection{اعداد ترتیبی جانشین و حدی}

در چند فصل آینده بسیار از اعداد ترتیبی استفاده خواهیم کرد. بنابراین معرفی
چند مفهوم ساده سودمند خواهد بود.

\begin{defn}
برای هر عدد ترتیبی~$\alpha$، \emph{جانشین} آن
$\ordsucc{\alpha}=\alpha\cup\{\alpha\}$ است. می‌گوییم~$\alpha$ یک عدد
ترتیبی \emph{جانشین} است اگر برای یک عدد ترتیبی~$\beta$،
$\ordsucc{\beta}=\alpha$. می‌گوییم~$\alpha$ یک عدد ترتیبی \emph{حدی}
است اگر و تنها اگر نه تهی باشد و نه عدد ترتیبی جانشین.
\end{defn}
\noindent
نتیجهٔ زیر نشان می‌دهد که این مفهوم درستی از \emph{جانشین} است:

\begin{prop}
برای هر عدد ترتیبی~$\alpha$:
\begin{enumerate}
	\item $\alpha \in \ordsucc{\alpha}$؛
	\item $\ordsucc{\alpha}$ یک عدد ترتیبی است؛
	\item هیچ عدد ترتیبی~$\beta$ وجود ندارد به‌طوری‌که
	$\alpha \in \beta \in \ordsucc{\alpha}$.
	\end{enumerate}
\end{prop}

\begin{proof}
بدیهی است که~$\alpha\in\alpha\cup\{\alpha\}=\ordsucc{\alpha}$.
همچنین~$\ordsucc{\alpha}$ مجموعه‌ای متعدی از اعداد ترتیبی است و بنابراین
بنا بر~\olref[basic]{corordtransitiveord} یک عدد ترتیبی است. و ممکن نیست
$\alpha\in\beta\in\ordsucc{\alpha}$، زیرا در آن صورت یا
$\beta\in\alpha$ یا~$\beta=\alpha$، برخلاف~\olref[basic]{ordordered}.
\end{proof}
\noindent
این نتیجه گونه‌ای از اثبات با استقرای ترامتناهی را نیز مجاز می‌سازد:

\begin{thm}[استقرای ترامتناهی ساده]\ollabel{simpletransrecursion}
فرض کنید~$\phi(x)$ فرمولی باشد به‌طوری‌که:
\begin{enumerate}
	\item $\phi(\emptyset)$؛ و
	\item برای هر عدد ترتیبی~$\alpha$، اگر~$\phi(\alpha)$، آنگاه
	$\phi(\ordsucc{\alpha})$؛ و
	\item اگر~$\alpha$ عدد ترتیبی حدی و
	$(\forall \beta \in \alpha)\phi(\beta)$ باشد، آنگاه~$\phi(\alpha)$.
\end{enumerate}
آنگاه~$\forall \alpha \phi(\alpha)$.
\end{thm}

\begin{proof}
عکس نقیض را اثبات می‌کنیم. پس فرض کنید عدد ترتیبی‌ای وجود دارد که
$\lnot\phi$ را ارضا می‌کند، و بگذارید~$\gamma$ کمینهٔ چنین اعداد ترتیبی
باشد. اگر~$\gamma=\emptyset$، این با بند نخست تناقض دارد. اگر برای یک عدد
ترتیبی~$\alpha$، $\gamma=\ordsucc{\alpha}$، کمینگی~$\gamma$ نتیجه
می‌دهد~$\phi(\alpha)$؛ سپس بند دوم~$\phi(\gamma)$ را به دست می‌دهد،
که باز تناقض است. و اگر~$\gamma$ عدد ترتیبی حدی باشد، کمینگی آن نتیجه
می‌دهد~$(\forall \beta \in \gamma)\phi(\beta)$؛ سپس بند سوم
$\phi(\gamma)$ را به دست می‌دهد، که تناقض است. این سه حالت بنا بر تعریف
همهٔ امکان‌ها را دربر می‌گیرند؛ پس هیچ عدد ترتیبی ناقض~$\phi$ وجود ندارد.
\end{proof}
\noindent
آخرین نمادگذاری زیر بعداً سودمند خواهد بود:

\begin{defn}\ollabel{defsupstrict}
اگر~$X$ مجموعه‌ای از اعداد ترتیبی باشد، آنگاه
$\supstrict(X)=\bigcup_{\alpha \in X}\ordsucc{\alpha}$.
\end{defn}
\noindent
در اینجا «lsub» مخفف «کمترین کران بالای اکید» است.\footnote{برخی کتاب‌ها
برای این مفهوم از «$\text{sup}(X)$» استفاده می‌کنند. اما کتاب‌های دیگری
«$\text{sup}(X)$» را برای کمترین کران بالای \emph{غیراکید}، یعنی صرفاً
$\bigcup X$، به کار می‌برند. اگر~$X$ بزرگ‌ترین عضوی مانند~$\alpha$ داشته
باشد، این دو مفهوم از هم جدا می‌شوند: کمترین کران بالای \emph{اکید}
$\ordsucc{\alpha}$ است، درحالی‌که کمترین کران بالای \emph{غیراکید}
صرفاً~$\alpha$ است.} نتیجهٔ زیر این مطلب را توضیح می‌دهد:

\begin{prop}
اگر~$X$ مجموعه‌ای از اعداد ترتیبی باشد، $\supstrict(X)$ کمترین عدد ترتیبی
بزرگ‌تر از هر عدد ترتیبی عضو~$X$ است.
\end{prop}

\begin{proof}
بگذارید~$Y=\Setabs{\ordsucc{\alpha}}{\alpha \in X}$، به‌طوری‌که
$\supstrict(X)=\bigcup Y$. چون اعداد ترتیبی متعدی‌اند و هر عضو یک عدد
ترتیبی خود عدد ترتیبی است، $\supstrict(X)$ مجموعه‌ای متعدی از اعداد
ترتیبی است و بنابراین بنا بر~\olref[basic]{corordtransitiveord} یک عدد
ترتیبی است.

اگر~$\alpha\in X$، آنگاه~$\ordsucc{\alpha}\in Y$، پس
$\ordsucc{\alpha}\subseteq\bigcup Y=\supstrict(X)$ و بنابراین
$\alpha\in\supstrict(X)$. پس~$\supstrict(X)$ از هر عدد ترتیبی عضو~$X$
اکیداً بزرگ‌تر است.

برعکس، اگر~$\alpha\in\supstrict(X)$، آنگاه برای یک~$\beta\in X$،
$\alpha\in\ordsucc{\beta}\in Y$ و در نتیجه~$\alpha\leq\beta$.

اکنون فرض کنید~$\gamma$ عدد ترتیبی دلخواهی باشد که از هر عدد ترتیبی
عضو~$X$ اکیداً بزرگ‌تر است. برای هر~$\alpha\in X$، از
$\alpha\in\gamma$ نتیجه می‌شود~$\ordsucc{\alpha}\subseteq\gamma$؛ پس
$\supstrict(X)=\bigcup_{\alpha \in X}\ordsucc{\alpha}\subseteq\gamma$.
بنابراین~$\supstrict(X)\leq\gamma$، و در نتیجه~$\supstrict(X)$
\emph{کمترین} کران بالای اکید~$X$ است.
\end{proof}

\end{document}
